%============================================================
% CH1 · P1 · PROBLEM 2 — MAXIMAL NETWORK RANK (LeetCode 1615)
%============================================================
\subsection{Problem 1.1.2: Maximal Network Rank}
\problemlabel{1.1.2}{ssec:ch1-network-rank}

\problemcard{1.1.2}{Maximal Network Rank}{LeetCode 1615}{https://leetcode.com/problems/maximal-network-rank/}{\exmedium\ \lc{Medium}}{Degrees of a pair + inclusion--exclusion on the shared edge}

\subsubsection*{1. Problem statement}

There are $n$ cities and a list of bidirectional roads; \texttt{roads[i]} $= [a_i, b_i]$ connects cities $a_i$ and $b_i$. The \textbf{network rank} of two different cities is the number of roads directly connected to \emph{either} city; a road connected to both cities is counted only once. The \textbf{maximal network rank} is the largest network rank over all pairs of different cities. Return it.

\textbf{Constraints:} $2 \le n \le 100$, \; $0 \le |\texttt{roads}| \le n(n-1)/2$, \; $0 \le a_i, b_i \le n-1$, \; $a_i \ne b_i$, \; each pair of cities has at most one road.

\begin{table}[H]
\centering\small
\begin{tabular}{@{}clcl@{}}
\toprule
$n$ & \textbf{roads} & \textbf{Output} & \textbf{A best pair} \\
\midrule
4 & $[[0,1],[0,3],[1,2],[1,3]]$ & 4 & cities 0 and 1 (road $0{-}1$ counted once) \\
5 & $[[0,1],[0,3],[1,2],[1,3],[2,3],[2,4]]$ & 5 & cities 1 and 2 \\
8 & $[[0,1],[1,2],[2,3],[2,4],[5,6],[5,7]]$ & 5 & cities 2 and 5 (in different parts!) \\
\bottomrule
\end{tabular}
\end{table}

\begin{diagrambox}
\begin{tikzpicture}[x=1cm, y=1cm]
    \begin{scope}
        \node[gvsel] (a0) at (0,0.7) {0}; \node[gvsel] (a1) at (1.6,0.7) {1};
        \node[gv] (a2) at (1.6,-0.7) {2}; \node[gv] (a3) at (0,-0.7) {3};
        \draw[gpath] (a0) -- (a1); \draw[gpath] (a0) -- (a3); \draw[gpath] (a1) -- (a2); \draw[gpath] (a1) -- (a3);
        \node[font=\scriptsize\bfseries] at (0.8,-1.5) {Example 1: rank$(0,1) = 4$};
    \end{scope}
    \begin{scope}[xshift=3.6cm]
        \node[gv] (b0) at (0,0) {0}; \node[gvsel] (b1) at (1.3,0.75) {1}; \node[gv] (b3) at (1.3,-0.75) {3};
        \node[gvsel] (b2) at (2.6,0) {2}; \node[gv] (b4) at (3.9,0) {4};
        \draw[gpath] (b0) -- (b1); \draw[ge] (b0) -- (b3); \draw[gpath] (b1) -- (b3);
        \draw[gpath] (b1) -- (b2); \draw[gpath] (b2) -- (b3); \draw[gpath] (b2) -- (b4);
        \node[font=\scriptsize\bfseries] at (1.95,-1.5) {Example 2: rank$(1,2) = 5$};
    \end{scope}
    \begin{scope}[xshift=9.2cm]
        \node[gv] (c0) at (-1.1,0.75) {0}; \node[gv] (c1) at (0,0.75) {1}; \node[gvsel] (c2) at (1.1,0.75) {2};
        \node[gv] (c3) at (0.45,-0.3) {3}; \node[gv] (c4) at (1.75,-0.3) {4};
        \node[gvsel] (c5) at (3.2,0.75) {5}; \node[gv] (c6) at (2.7,-0.3) {6}; \node[gv] (c7) at (3.7,-0.3) {7};
        \draw[ge] (c0) -- (c1); \draw[gpath] (c1) -- (c2); \draw[gpath] (c2) -- (c3); \draw[gpath] (c2) -- (c4);
        \draw[gpath] (c5) -- (c6); \draw[gpath] (c5) -- (c7);
        \node[font=\scriptsize\bfseries] at (1.25,-1.5) {Example 3: rank$(2,5) = 5$};
    \end{scope}
\end{tikzpicture}
\figcap{The three official examples. Purple: a best pair; orange: the roads it touches. In Example 2 the road $1{-}2$ touches both cities but counts once.}
\end{diagrambox}

\subsubsection*{2. The graph model}
\begin{modelbox}[Graph Model --- Maximal Network Rank]
\renewcommand{\arraystretch}{1.35}
\begin{tabularx}{\linewidth}{@{}>{\bfseries}l X@{}}
Vertices & the $n$ cities, $0..n-1$ \\
Edges & one undirected edge per road; a simple graph (no loops, no parallel roads); need not be connected \\
Weights & none \\
Question & $\displaystyle\max_{u \ne v}\ \big|E(u) \cup E(v)\big| = \max_{u \ne v}\ \deg(u) + \deg(v) - [\,uv \in E\,]$, where $E(x)$ is the set of edges at $x$ \\
Size & $|V| \le 100$, $|E| \le 4950$; $\binom{100}{2} = 4950$ pairs: even $\bigO{|V|^2}$ is tiny \\
\end{tabularx}
\end{modelbox}

\begin{diagrambox}
\begin{tikzpicture}[x=1.1cm, y=1cm]
    \begin{scope}
        \node[gvsel] (u) at (0,0) {$u$}; \node[gvsel] (v) at (2,0) {$v$};
        \node[gv, minimum size=5mm] (x1) at (-1,0.8) {}; \node[gv, minimum size=5mm] (x2) at (-1,-0.8) {};
        \node[gv, minimum size=5mm] (y1) at (3,0.8) {}; \node[gv, minimum size=5mm] (y2) at (3,-0.8) {};
        \draw[draw=primaryblue, line width=1.3pt] (u) -- (x1); \draw[draw=primaryblue, line width=1.3pt] (u) -- (x2);
        \draw[draw=accentorange, line width=1.3pt] (v) -- (y1); \draw[draw=accentorange, line width=1.3pt] (v) -- (y2);
        \node[font=\scriptsize\bfseries, align=center] at (1,-1.45) {not adjacent:\\ rank $= 2 + 2 = 4$};
    \end{scope}
    \begin{scope}[xshift=6.2cm]
        \node[gvsel] (u) at (0,0) {$u$}; \node[gvsel] (v) at (2,0) {$v$};
        \node[gv, minimum size=5mm] (x1) at (-1,0.8) {}; \node[gv, minimum size=5mm] (x2) at (-1,-0.8) {};
        \node[gv, minimum size=5mm] (y1) at (3,0.8) {}; \node[gv, minimum size=5mm] (y2) at (3,-0.8) {};
        \draw[draw=primaryblue, line width=1.3pt] (u) -- (x1); \draw[draw=primaryblue, line width=1.3pt] (u) -- (x2);
        \draw[draw=accentorange, line width=1.3pt] (v) -- (y1); \draw[draw=accentorange, line width=1.3pt] (v) -- (y2);
        \draw[gpath, draw=softpurple] (u) -- (v);
        \node[font=\scriptsize\bfseries, align=center] at (1,-1.45) {adjacent: $\deg = 3, 3$\\ rank $= 3 + 3 - 1 = 5$};
    \end{scope}
\end{tikzpicture}
\figcap{The rank of a pair is a pure degree formula. The only non-local ingredient is one bit: is $uv$ an edge?}
\end{diagrambox}

\subsubsection*{3. Key observations}

\begin{itemize}[leftmargin=*]
    \item \textbf{Inclusion--exclusion.} The roads touching $u$ or $v$ are the $\deg(u)$ roads at $u$ plus the $\deg(v)$ roads at $v$, minus the ones counted twice. Only the road $uv$ itself can touch both (the graph is simple), so $\text{rank}(u,v) = \deg(u) + \deg(v) - [uv \in E]$.
    \item \textbf{Connectivity is irrelevant.} The two cities need not be connected at all (Example 3). Nothing here needs a traversal --- Pattern 1.
    \item \textbf{``Take the two largest degrees'' is not enough.} The two largest degrees might belong to adjacent cities while another pair with the same degree sum is not adjacent. On a 4-cycle every degree is 2; the first two vertices $0,1$ are adjacent (rank 3), but $0,2$ are not (rank 4).
    \item \textbf{But only top degrees matter.} The correction is at most 1, so an optimal pair always uses the largest degree and, if that is unique, the second largest. This gives an $\bigO{|V|+|E|}$ algorithm (Approach 3), needed when $n$ is large.
\end{itemize}

\begin{diagrambox}
\begin{tikzpicture}[x=1cm, y=1cm]
    \begin{scope}
        \node[gvsel] (p0) at (0,1.2) {0}; \node[gvsel] (p1) at (1.4,1.2) {1};
        \node[gv] (p2) at (1.4,-0.2) {2}; \node[gv] (p3) at (0,-0.2) {3};
        \draw[gpath, draw=softpurple] (p0) -- (p1); \draw[ge] (p1) -- (p2); \draw[ge] (p2) -- (p3); \draw[ge] (p3) -- (p0);
        \node[font=\scriptsize\bfseries, align=center] at (0.7,-1.1) {greedy: first two of degree 2\\ $2 + 2 - 1 = 3$ \no};
    \end{scope}
    \begin{scope}[xshift=5cm]
        \node[gvsel] (q0) at (0,1.2) {0}; \node[gv] (q1) at (1.4,1.2) {1};
        \node[gvsel] (q2) at (1.4,-0.2) {2}; \node[gv] (q3) at (0,-0.2) {3};
        \draw[gpath] (q0) -- (q1); \draw[gpath] (q1) -- (q2); \draw[gpath] (q2) -- (q3); \draw[gpath] (q3) -- (q0);
        \node[font=\scriptsize\bfseries, align=center] at (0.7,-1.1) {opposite corners, not adjacent\\ $2 + 2 - 0 = 4$ \yes};
    \end{scope}
    \node[note, anchor=west, text width=3.6cm, align=left] at (8.1,0.5) {Ties among the top degrees must be examined: one non-adjacent pair is enough.};
\end{tikzpicture}
\figcap{The 4-cycle: picking ``the two largest degrees'' blindly gives 3; the answer is 4.}
\end{diagrambox}

\subsubsection*{4. Approaches}

\begin{bruteforcebox}
\textbf{Idea.} Follow the definition literally: for every pair $(u,v)$, scan all roads and count those touching $u$ or $v$.
\begin{lstlisting}[style=cpp]
#include <bits/stdc++.h>
using namespace std;

class Solution {
public:
    int maximalNetworkRank(int n, vector<vector<int>>& roads) {
        int best = 0;
        for (int u = 0; u < n; u++)
            for (int v = u + 1; v < n; v++) {
                int cnt = 0;                               // roads touching u or v
                for (auto& r : roads)
                    if (r[0] == u || r[1] == u || r[0] == v || r[1] == v) cnt++;
                best = max(best, cnt);
            }
        return best;
    }
};
\end{lstlisting}
A road touching both cities satisfies the condition once and is counted once --- correct by definition, no inclusion--exclusion needed.
\begin{complexitybox}
\textbf{Time:} $\bigO{n^2 \cdot |E|}$ --- for $n = 100$, $|E| = 4950$ about $2.5 \cdot 10^7$ checks: \ac{} here, hopeless for $n = 10^5$. \textbf{Space:} $\bigO{1}$.
\end{complexitybox}
\end{bruteforcebox}

\begin{algorithmbox}
\textbf{Idea.} One pass computes all degrees and fills an adjacency matrix; then every pair costs $\bigO{1}$ with the formula $\deg(u) + \deg(v) - [uv \in E]$.
\begin{lstlisting}[style=cpp]
#include <bits/stdc++.h>
using namespace std;

class Solution {
public:
    int maximalNetworkRank(int n, vector<vector<int>>& roads) {
        vector<int> deg(n, 0);
        vector<vector<bool>> adj(n, vector<bool>(n, false));
        for (auto& r : roads) {
            deg[r[0]]++; deg[r[1]]++;                  // Pattern 1 counting
            adj[r[0]][r[1]] = adj[r[1]][r[0]] = true;  // remember the edge (both directions)
        }
        int best = 0;
        for (int u = 0; u < n; u++)
            for (int v = u + 1; v < n; v++)
                best = max(best, deg[u] + deg[v] - (adj[u][v] ? 1 : 0));
        return best;
    }
};
\end{lstlisting}
\begin{complexitybox}
\textbf{Time:} $\bigO{n^2 + |E|}$ --- about $5\,000$ pair evaluations. \textbf{Space:} $\bigO{n^2}$ bits for the matrix ($10^4$ here). \textbf{Verdict:} \ac{}.
\end{complexitybox}
\end{algorithmbox}

\begin{invariantbox}[Why the formula is the network rank]
Let $E(x)$ be the set of roads at city $x$, so $|E(x)| = \deg(x)$. The network rank is $|E(u) \cup E(v)| = |E(u)| + |E(v)| - |E(u) \cap E(v)|$. A road in the intersection has both $u$ and $v$ as endpoints, so it is the road $uv$; since parallel roads are forbidden there is at most one. Hence $|E(u) \cap E(v)| = [uv \in E]$ and the loop evaluates the exact rank of every pair; the maximum over all pairs is the answer. The degrees are correct by the counting invariant of Pattern 1. \hfill$\blacksquare$
\end{invariantbox}

\begin{alternativebox}[{Approach 3: Only the Top Degrees --- $\bigO{|V|+|E|}$}]
Let $d_1$ be the largest degree and $c_1$ the number of cities having it; let $d_2 < d_1$ be the largest smaller degree and $c_2$ its count.
\begin{itemize}[leftmargin=*]
    \item \textbf{$c_1 \ge 2$.} Every pair has rank $\le 2d_1$, reached exactly by a \emph{non-adjacent} pair of top cities. If all $\binom{c_1}{2}$ top pairs are adjacent, two top cities still give $2d_1 - 1$, and any other pair gives at most $d_1 + d_2 \le 2d_1 - 1$. So: count the roads with both ends of degree $d_1$; answer $2d_1$ if this count is $< \binom{c_1}{2}$, else $2d_1 - 1$.
    \item \textbf{$c_1 = 1$} (one top city $x$). Pairs without $x$ give $\le 2d_2 \le d_1 + d_2 - 1$; pairs $(x, w)$ with $\deg w < d_2$ give $\le d_1 + d_2 - 1$. So the answer is $d_1 + d_2$ if $x$ misses at least one city of degree $d_2$, else $d_1 + d_2 - 1$: count the roads joining $x$ to a degree-$d_2$ city and compare with $c_2$.
\end{itemize}
\begin{lstlisting}[style=cpp]
#include <bits/stdc++.h>
using namespace std;

class Solution {
public:
    int maximalNetworkRank(int n, vector<vector<int>>& roads) {
        vector<int> deg(n, 0);
        for (auto& r : roads) { deg[r[0]]++; deg[r[1]]++; }
        int d1 = -1, d2 = -1;                          // top two distinct degrees
        for (int v = 0; v < n; v++) {
            if (deg[v] > d1) { d2 = d1; d1 = deg[v]; }
            else if (deg[v] < d1 && deg[v] > d2) d2 = deg[v];
        }
        long long c1 = 0, c2 = 0;
        for (int v = 0; v < n; v++) {
            if (deg[v] == d1) c1++;
            else if (deg[v] == d2) c2++;
        }
        if (c1 >= 2) {                                 // several top cities
            long long inside = 0;                      // roads between two top cities
            for (auto& r : roads) if (deg[r[0]] == d1 && deg[r[1]] == d1) inside++;
            return inside < c1 * (c1 - 1) / 2 ? 2 * d1 : 2 * d1 - 1;
        }
        long long cross = 0;                           // roads: top city -- d2-city
        for (auto& r : roads)
            if ((deg[r[0]] == d1 && deg[r[1]] == d2) ||
                (deg[r[0]] == d2 && deg[r[1]] == d1)) cross++;
        return cross < c2 ? d1 + d2 : d1 + d2 - 1;
    }
};
\end{lstlisting}
When $c_1 = 1$, $d_2$ always exists because $n \ge 2$. The counting trick replaces ``check every pair'' by ``count edges inside a group and compare with the number of pairs in it''.

\begin{diagrambox}
\begin{tikzpicture}[x=1cm, y=1cm]
    \begin{scope}
        \node[gv] (b0) at (0,0) {0}; \node[gvsel] (b1) at (1.3,0.75) {1}; \node[gvsel] (b3) at (1.3,-0.75) {3};
        \node[gvsel] (b2) at (2.6,0) {2}; \node[gv] (b4) at (3.9,0) {4};
        \draw[ge] (b0) -- (b1); \draw[ge] (b0) -- (b3); \draw[gpath, draw=softpurple] (b1) -- (b3);
        \draw[gpath, draw=softpurple] (b1) -- (b2); \draw[gpath, draw=softpurple] (b2) -- (b3); \draw[ge] (b2) -- (b4);
        \node[font=\scriptsize, align=center] at (1.95,-1.55) {Example 2: $d_1 = 3$, $c_1 = 3$; 3 roads inside, $\binom32 = 3$\\ all top pairs adjacent $\Rightarrow 2\cdot3-1 = \mathbf{5}$};
    \end{scope}
    \begin{scope}[xshift=7.2cm]
        \node[gv] (c0) at (-1.1,0.75) {0}; \node[gvseen] (c1) at (0,0.75) {1}; \node[gvsel] (c2) at (1.1,0.75) {2};
        \node[gv] (c3) at (0.45,-0.3) {3}; \node[gv] (c4) at (1.75,-0.3) {4};
        \node[gvseen] (c5) at (3.2,0.75) {5}; \node[gv] (c6) at (2.7,-0.3) {6}; \node[gv] (c7) at (3.7,-0.3) {7};
        \draw[ge] (c0) -- (c1); \draw[gpath, draw=softpurple] (c1) -- (c2); \draw[ge] (c2) -- (c3); \draw[ge] (c2) -- (c4);
        \draw[ge] (c5) -- (c6); \draw[ge] (c5) -- (c7);
        \node[font=\scriptsize, align=center] at (1.25,-1.55) {Example 3: top city 2 ($d_1=3$); $d_2 = 2$: cities 1, 5\\ 1 cross road $< c_2 = 2 \Rightarrow 3+2 = \mathbf{5}$};
    \end{scope}
\end{tikzpicture}
\figcap{Approach 3 on Examples 2 and 3. Purple vertices: degree $d_1$; yellow: degree $d_2$; purple edges: the roads that are counted.}
\end{diagrambox}
\end{alternativebox}

\begin{tricksbox}
\begin{itemize}[leftmargin=*]
    \item \textbf{Adjacency test without a matrix:} for large $n$, store each road as the key $\min \cdot n + \max$ in a hash set (Pattern 1 template 2), or sort each adjacency list and binary-search.
    \item \textbf{\texttt{vector<vector<bool>>}} packs bits; for speed in tight loops a flat \texttt{vector<char>} of size $n^2$ with index \texttt{u*n+v} is faster.
    \item \textbf{Upper bound pruning:} sort cities by degree; once $\deg(u) + \deg(v) \le \text{best}$ the inner loop can stop. With the correction at most 1, the search ends almost immediately --- that is Approach 3 in disguise.
\end{itemize}
\end{tricksbox}

\subsubsection*{5. Trace}

\begin{tracebox}[{Example 2: $n = 5$, roads $01, 03, 12, 13, 23, 24$}]
\centering
\begin{tikzpicture}[x=1cm, y=1cm]
    \begin{scope}
        \node[gv] (b0) at (0,0) {0}; \node[gv] (b1) at (1.3,0.8) {1}; \node[gv] (b3) at (1.3,-0.8) {3};
        \node[gv] (b2) at (2.6,0) {2}; \node[gv] (b4) at (3.9,0) {4};
        \draw[ge] (b0) -- (b1); \draw[ge] (b0) -- (b3); \draw[ge] (b1) -- (b3);
        \draw[ge] (b1) -- (b2); \draw[ge] (b2) -- (b3); \draw[ge] (b2) -- (b4);
        \node[gdist, left=1mm of b0] {2}; \node[gdist, above=1mm of b1] {3}; \node[gdist, below=1mm of b3] {3};
        \node[gdist, below right=0.3mm of b2] {3}; \node[gdist, above=1mm of b4] {1};
        \node[font=\scriptsize\bfseries] at (1.95,-1.75) {step 1: degrees $2, 3, 3, 3, 1$ (sum $12 = 2\cdot 6$)};
    \end{scope}
    \begin{scope}[xshift=6.6cm, yshift=1.3cm]
        \foreach \i in {0,...,4}{
            \node[idx] at (0.45+\i*0.8,0.55) {\i};
            \node[idx] at (-0.2,-\i*0.6) {\i};
        }
        \foreach \r/\row in {0/{-,4,5,4,3}, 1/{4,-,5,5,4}, 2/{5,5,-,5,3}, 3/{4,5,5,-,4}, 4/{3,4,3,4,-}}{
            \foreach \x [count=\c from 0] in \row {
                \ifnum\c=\r
                    \node[emptycell, minimum width=8mm, minimum height=6mm] at (0.45+\c*0.8,-\r*0.6) {};
                \else
                    \node[cell, minimum width=8mm, minimum height=6mm] at (0.45+\c*0.8,-\r*0.6) {\x};
                \fi
            }
        }
        \foreach \r/\c in {0/2, 1/2, 1/3, 2/3}{
            \node[anscell, minimum width=8mm, minimum height=6mm] at (0.45+\c*0.8,-\r*0.6) {5};
        }
        \node[font=\scriptsize\bfseries] at (2.05,-3.05) {step 2: rank$(u,v)$ for every pair};
    \end{scope}
\end{tikzpicture}
\figcap{Left: degrees from one pass. Right: the rank table (upper triangle highlighted where the maximum 5 is reached: pairs $(0,2), (1,2), (1,3), (2,3)$).}

\medskip
\small
\begin{tabular}{@{}lccccl@{}}
\toprule
\textbf{pair} & $\deg u$ & $\deg v$ & adjacent? & \textbf{rank} & \textbf{best so far} \\
\midrule
$(0,1)$ & 2 & 3 & yes & 4 & 4 \\
$(0,2)$ & 2 & 3 & no  & 5 & 5 \\
$(0,3)$ & 2 & 3 & yes & 4 & 5 \\
$(0,4)$ & 2 & 1 & no  & 3 & 5 \\
$(1,2)$ & 3 & 3 & yes & 5 & 5 \\
$(1,3)$ & 3 & 3 & yes & 5 & 5 \\
$(1,4)$ & 3 & 1 & no  & 4 & 5 \\
$(2,3)$ & 3 & 3 & yes & 5 & 5 \\
$(2,4)$ & 3 & 1 & yes & 3 & 5 \\
$(3,4)$ & 3 & 1 & no  & 4 & 5 \\
\bottomrule
\end{tabular}
\end{tracebox}

\begin{tracebox}[{Example 2: the winning pairs on the map}]
\centering
\begin{tikzpicture}[x=1cm, y=1cm]
    \begin{scope}
        \node[gvsel] (b0) at (0,0) {0}; \node[gv] (b1) at (1.3,0.8) {1}; \node[gv] (b3) at (1.3,-0.8) {3};
        \node[gvsel] (b2) at (2.6,0) {2}; \node[gv] (b4) at (3.9,0) {4};
        \draw[gpath] (b0) -- (b1); \draw[gpath] (b0) -- (b3); \draw[ge] (b1) -- (b3);
        \draw[gpath] (b1) -- (b2); \draw[gpath] (b2) -- (b3); \draw[gpath] (b2) -- (b4);
        \node[font=\scriptsize\bfseries, align=center] at (1.95,-1.6) {$(0,2)$: $2 + 3 - 0 = 5$\\ not adjacent, nothing subtracted};
    \end{scope}
    \begin{scope}[xshift=6.6cm]
        \node[gv] (b0) at (0,0) {0}; \node[gvsel] (b1) at (1.3,0.8) {1}; \node[gv] (b3) at (1.3,-0.8) {3};
        \node[gvsel] (b2) at (2.6,0) {2}; \node[gv] (b4) at (3.9,0) {4};
        \draw[gpath] (b0) -- (b1); \draw[ge] (b0) -- (b3); \draw[gpath] (b1) -- (b3);
        \draw[gpath, draw=softpurple] (b1) -- (b2); \draw[gpath] (b2) -- (b3); \draw[gpath] (b2) -- (b4);
        \node[font=\scriptsize\bfseries, align=center] at (1.95,-1.6) {$(1,2)$: $3 + 3 - 1 = 5$\\ road $1{-}2$ (purple) counted once};
    \end{scope}
\end{tikzpicture}
\figcap{Two different reasons for the same maximum: a smaller degree sum with no overlap, or a larger sum minus the shared road.}
\end{tracebox}

\subsubsection*{6. Test suite}

\begin{testbox}
\centering\small
\begin{tabular}{@{}cllc@{}}
\toprule
$n$ & \textbf{roads} & \textbf{Why it is interesting} & \textbf{Expected} \\
\midrule
2 & $[\,]$ & no roads at all & 0 \\
2 & $[[0,1]]$ & one road touching both cities: counted once & 1 \\
4 & $[[0,1],[1,2],[2,3],[3,0]]$ & 4-cycle: the tie trap of the greedy & 4 \\
4 & all 6 roads ($K_4$) & every top pair adjacent: $2\cdot3-1$ & 5 \\
5 & $[[0,1],[0,2],[0,3],[0,4]]$ & star: center + any leaf, $4 + 1 - 1$ & 4 \\
8 & official Example 3 & best pair in different components & 5 \\
100 & all $4950$ roads ($K_{100}$) & largest input: $99 + 99 - 1$ & 197 \\
\bottomrule
\end{tabular}
\end{testbox}

\subsubsection*{7. Common mistakes}

\begin{warningbox}[Mistake 1 --- forgetting the shared road]
\begin{lstlisting}[style=snippet]
best = max(best, deg[u] + deg[v]);                          // WRONG: Example 1 gives 5
best = max(best, deg[u] + deg[v] - (adj[u][v] ? 1 : 0));    // RIGHT: 4
\end{lstlisting}
\end{warningbox}

\begin{warningbox}[Mistake 2 --- taking the two largest degrees without looking at ties]
Sorting degrees and returning $d_{(1)} + d_{(2)} - [\text{those two cities adjacent}]$ fails on the 4-cycle (3 instead of 4): another pair with the same degree sum may be non-adjacent. Either check all pairs, or count edges inside the top group (Approach 3).
\end{warningbox}

\begin{warningbox}[Mistake 3 --- assuming the graph is connected]
Running a BFS from city 0 and only considering reachable cities misses Example 3. The rank is defined for \emph{every} pair.
\end{warningbox}

\begin{warningbox}[Mistake 4 --- filling the matrix in one direction only]
\texttt{adj[a][b] = true} alone makes \texttt{adj[u][v]} with $u < v$ false whenever the road was given as $[v, u]$. Set both entries (or normalise each pair to $(\min, \max)$).
\end{warningbox}

\subsubsection*{8. Takeaways}

\begin{itemize}[leftmargin=*]
    \item A pair question about ``edges touching $u$ or $v$'' is inclusion--exclusion: $\deg u + \deg v$ minus the shared edges.
    \item Degrees answer the local part; a single adjacency bit fixes the overlap.
    \item When the correction is bounded (here by 1), only the top one or two degree classes can win --- the key to linear time.
    \item Draw the definition on a tiny example (the 4-cycle) before trusting a greedy.
\end{itemize}

\begin{variantsbox}
\variant{Maximum Star Sum of a Graph}{\varin{https://leetcode.com/problems/maximum-star-sum-of-a-graph/}}
  {vertices carry values; choose a center and at most $k$ of its neighbours to maximise the sum.}
  {model: again a purely local question --- each center only looks at its own incident edges.}
  {algorithm: for every vertex keep the positive neighbour values, take the $k$ largest; $\bigO{|V|+|E|\log|E|}$.}
\variant{Maximum Total Importance of Roads}{\varin{https://leetcode.com/problems/maximum-total-importance-of-roads/}}
  {assign the values $1..n$ to the cities to maximise $\sum_{\text{roads}} (\text{val}(a)+\text{val}(b))$.}
  {model: the sum over edges equals $\sum_v \text{val}(v)\deg(v)$ (Tool 4 of Pattern 1).}
  {algorithm: sort degrees; the largest value goes to the largest degree (rearrangement inequality).}
\variant{Network rank for $n, |E| \le 2\cdot10^5$}{\varnone}
  {the same question on a big sparse graph.}
  {model: unchanged; the $n^2$ pair loop and the matrix no longer fit.}
  {algorithm: Approach 3 --- only the top one or two degree classes matter; count edges inside them; $\bigO{|V|+|E|}$.}
\variant{Remove two vertices from a tree}{\varnone}
  {delete two vertices $u \ne v$ (with their edges) from a tree; maximise the number of remaining components.}
  {model: removing $u$ splits the tree into $\deg(u)$ parts; removing $v$ afterwards adds $\deg(v) - 1 - [uv \in E]$.}
  {algorithm: maximise $\deg u + \deg v - 1 - [uv \in E]$ --- exactly this problem's formula shifted by $-1$, solved by Approach 3.}
\end{variantsbox}
